\documentclass[a4paper,twoside]{article}

\usepackage{morewrites}

\usepackage[svgnames,dvipsnames,x11names]{xcolor}
\usepackage{graphicx}
\usepackage[english]{babel}
\usepackage[colorlinks=true, linkcolor=NavyBlue, urlcolor=NavyBlue, citecolor=NavyBlue]{hyperref}
\usepackage[a4paper]{geometry}
\usepackage{ts-fonts}
\usepackage{amsmath}
\usepackage{float}
\usepackage{seqsplit}
\usepackage{ts-callouts}
\usepackage{ts-code}

\usepackage{lastpage}
\usepackage{refcount}
\makeatletter
\def\ps@plain{
  \let\@oddhead\@empty
  \let\@evenhead\@empty
  \def\@oddfoot{
    \hfil
    \ifnum\getpagerefnumber{LastPage}>1\relax
      \thepage
    \fi
    \hfil}
  \let\@evenfoot\@oddfoot
}
\makeatother

\title{Text Formatting}
\author{}\date{}

\begin{document}

\maketitle

\thispagestyle{plain}
Like vanilla Markdown, you can apply basic text formatting using a variety of delimiters. TeXSmith extends this with small capitals support.

\begin{code}{markdown}{}{}
\begin{Verbatim}[commandchars=\\\{\},breaklines, breakanywhere, commandchars=\\\{\}]
The quick brown fox jumps over the lazy dog. \PY{g+ge}{*(regular)*}

\PY{g+ge}{*The quick brown fox jumps over the lazy dog.*} \PY{g+ge}{*(italic)*}

\PY{g+gs}{**The quick brown fox jumps over the lazy dog.**} \PY{g+ge}{*(bold)*}

***The quick brown fox jumps over the lazy dog.*** \PY{g+ge}{*(bold italic)*}

\PY{g+gd}{\PYZti{}\PYZti{}The quick brown fox jumps over the lazy dog.\PYZti{}\PYZti{}} \PY{g+ge}{*(strikethrough)*}

\PY{g+gs}{\PYZus{}\PYZus{}The quick brown fox jumps over the lazy dog.\PYZus{}\PYZus{}} \PY{g+ge}{*(small capitals)*}
\end{Verbatim}
\end{code}
\begin{figure}[H]
    \centering

    \ifdefined\adjustbox

        \adjustbox{max width=\textwidth}{
            \includegraphics[width=\linewidth]{snippet-<HASH>.pdf}
        }

    \else

        \includegraphics[width=\linewidth]{snippet-<HASH>.pdf}

    \fi

\end{figure}
The \texttt{pymdownx.betterem} extension lets you stack delimiters for bold italic.

\section{Standalone bold paragraphs (\texttt{\textbackslash{}tslead})}\label{standalone-bold-paragraphs-texttt-textbackslash-tslead}

A paragraph whose only content is a short bold span (under 80 characters) is promoted to a lead-in pseudo-heading rather than rendered as a plain \texttt{\textbackslash{}textbf\{…\}}. TeXSmith emits \texttt{\textbackslash{}tslead\{…\}}, defined as:

\begin{code}{latex}{}{}
\begin{Verbatim}[commandchars=\\\{\},breaklines, breakanywhere, commandchars=\\\{\}]
\PY{k}{\PYZbs{}providecommand}\PY{n+nb}{\PYZob{}}\PY{k}{\PYZbs{}tslead}\PY{n+nb}{\PYZcb{}}[1]\PY{n+nb}{\PYZob{}}\PY{k}{\PYZbs{}par}\PY{k}{\PYZbs{}noindent}\PY{k}{\PYZbs{}textbf}\PY{n+nb}{\PYZob{}}\PYZsh{}1\PY{n+nb}{\PYZcb{}}\PY{k}{\PYZbs{}par}\PY{k}{\PYZbs{}nobreak}\PY{k}{\PYZbs{}smallskip}\PY{n+nb}{\PYZcb{}}
\end{Verbatim}
\end{code}
This guarantees a no-indent paragraph break and a small vertical breather, so the label looks identical regardless of what precedes it. Without this, \texttt{**Méthodologie**} wedged between two tables (\texttt{\textbackslash{}end\{center\}} ... \texttt{\textbackslash{}textbf\{…\}} ... \texttt{\textbackslash{}begin\{center\}}) would render flush left while the same construct after running prose would be indented by \texttt{babel-\allowbreak{}french}'s \texttt{\textbackslash{}parindent}. With \texttt{\textbackslash{}tslead}, both cases align.

\begin{code}{markdown}{}{}
\begin{Verbatim}[commandchars=\\\{\},breaklines, breakanywhere, commandchars=\\\{\}]
La synthèse récapitule, pilier par pilier, les forces et faiblesses…

\PY{g+gs}{**Sens critique**}

| Python | C |
|\PYZhy{}\PYZhy{}\PYZhy{}\PYZhy{}\PYZhy{}\PYZhy{}\PYZhy{}\PYZhy{}|\PYZhy{}\PYZhy{}\PYZhy{}|
| Faible | Fort |

\PY{g+gs}{**Méthodologie**}

| Python | C |
|\PYZhy{}\PYZhy{}\PYZhy{}\PYZhy{}\PYZhy{}\PYZhy{}\PYZhy{}\PYZhy{}|\PYZhy{}\PYZhy{}\PYZhy{}|
| Correcte | Forte |
\end{Verbatim}
\end{code}
The rule fires when the \texttt{<p>} contains exactly one \texttt{<strong>} child and nothing else (whitespace aside), and the bold's plain-text content is shorter than 80 characters. Bold spans inside running prose (\texttt{Some **bold** text.}), bold paragraphs over the threshold, and bold labels synthesised by other extensions (e.g. \texttt{tabbed-\allowbreak{}set} labels, which sit inside a \texttt{<div>}) keep their original \texttt{\textbackslash{}textbf\{…\}} rendering.

Override \texttt{\textbackslash{}tslead} in a custom preamble snippet to change the visual style --- for instance, to add a coloured rule, switch to small caps, or replace the \texttt{\textbackslash{}smallskip} with \texttt{\textbackslash{}medskip}:

\begin{code}{latex}{}{}
\begin{Verbatim}[commandchars=\\\{\},breaklines, breakanywhere, commandchars=\\\{\}]
\PY{k}{\PYZbs{}renewcommand}\PY{n+nb}{\PYZob{}}\PY{k}{\PYZbs{}tslead}\PY{n+nb}{\PYZcb{}}[1]\PY{n+nb}{\PYZob{}}\PY{k}{\PYZbs{}par}\PY{k}{\PYZbs{}noindent}\PY{k}{\PYZbs{}textsc}\PY{n+nb}{\PYZob{}}\PYZsh{}1\PY{n+nb}{\PYZcb{}}\PY{k}{\PYZbs{}par}\PY{k}{\PYZbs{}nobreak}\PY{k}{\PYZbs{}medskip}\PY{n+nb}{\PYZcb{}}
\end{Verbatim}
\end{code}
\begin{callout}[callout note]{Note}
In MkDocs, you need to specify how to render small capitals using a custom CSS:

\begin{code}{css}{}{}
\begin{Verbatim}[commandchars=\\\{\},breaklines, breakanywhere, commandchars=\\\{\}]
\PY{p}{.}\PY{n+nc}{texsmith\PYZhy{}smallcaps}\PY{+w}{ }\PY{p}{\PYZob{}}
\PY{+w}{    }\PY{k}{font\PYZhy{}variant}\PY{p}{:}\PY{+w}{ }\PY{k+kc}{small\PYZhy{}caps}\PY{p}{;}
\PY{+w}{    }\PY{k}{letter\PYZhy{}spacing}\PY{p}{:}\PY{+w}{ }\PY{l+m+mf}{0.04}\PY{k+kt}{em}\PY{p}{;}
\PY{p}{\PYZcb{}}
\end{Verbatim}
\end{code}
Then, include this CSS in your MkDocs configuration under \texttt{extra\_css}:

\begin{code}{yaml}{}{}
\begin{Verbatim}[commandchars=\\\{\},breaklines, breakanywhere, commandchars=\\\{\}]
\PY{n+nt}{extra\PYZus{}css}\PY{p}{:}
\PY{+w}{  }\PY{p+pIndicator}{\PYZhy{}}\PY{+w}{ }\PY{l+lScalar+lScalarPlain}{stylesheets/smallcaps.css}
\end{Verbatim}
\end{code}
\end{callout}

\end{document}
